% Historical source (removed by author request): references/_OLD_Übungsaufgaben/Übungsblatt 7 - Komplexe Zahlen/Blatt7.tex:66
\begin{ExerciseBlock}
\ExerciseHeading{4}{Darstellungswechsel}{complex_numbers_representation_conversion}
\begin{enumerate}[label=\alph*)]
	\item Bestimmen Sie für die folgenden komplexen Zahlen in kartesischer Darstellung die dazugehörigen polaren Darstellungen (trigonometrische und exponentielle Darstellung):
	\begin{enumerate}[label=\arabic*)]
		\item $z = 3 + 4i$
		\item $z = -1 + i$
	\end{enumerate}
	
	\item Bestimmen Sie für die folgenden komplexen Zahlen in polarer Darstellung die dazugehörige kartesische Darstellung:
	\begin{enumerate}[label=\arabic*)]
		\item $z = \sqrt{8} \; e^{i \frac{\pi}{4} }$
		\item $z = 2 \left( \cos\left( \pi \right) + i \sin\left( \pi \right) \right)$
	\end{enumerate}
\end{enumerate}

\end{ExerciseBlock}
\UseExerciseSolution{complex_numbers_representation_conversion}
